\documentclass[../main.tex]{subfiles}
\begin{document}
%==========================================TIÊU ĐỀ TẬP========================================
\begin{table}[h]
    \begin{adjustwidth}{-1cm}{}
        \begin{tabular}{>{\centering\arraybackslash}p{6cm}|>{\raggedright\arraybackslash}p{12.5cm}}
            \multirow{5}{6cm}
            {\\[-40pt]
                \flaregothic\fontsize{20pt}{20pt}\selectfont \centering
                \color{eptype}HISTOIRE\color{black}\\
                \freshbot\fontsize{72pt}{72pt}\selectfont
                \color{epnum}33\\[6pt]
                \cafeteria\fontsize{20pt}{20pt}\selectfont\color{epnum}(D-33)\color{black}
                \\[18pt]
                \color{black}
            }
             &
            % Dòng thứ nhất: tên tập tiếng Nhật
            {\fotskip\fontsize{18pt}{18pt}\selectfont \ruby{廃}{はい}\ruby{工}{こう}\ruby{場}{じょう}\ruby{跡}{あと}の\ruby{実}{じっ}\ruby{験}{けん}\ruby{室}{しつ}}
            \\[6pt]
            % Dòng thứ hai: tên tập tiếng Anh
             & {\cafeteria\fontsize{18pt}{18pt}\selectfont The Laboratory on Abandoned Ground}                                                     \\[4pt]
            % Dòng thứ ba: tên tập tiếng Việt
             & {\vnmsans\fontsize{18pt}{18pt}\selectfont Phòng thí nghiệm trên bãi đất hoang}                                                      \\[8pt]
            % Dòng thứ tư: Nhân vật xuất hiện
             & {\caviar{\fontsize{13pt}{13pt}\selectfont Nhân vật: \emp{suzuno2} \emp{kuon2} \emp{kaiho2} \emp{ryuuhou2} \emp{naomi2} \emp{game}}}
            \\[8pt]
        \end{tabular}
    \end{adjustwidth}
\end{table}
%=========================================TRUYỆN CHÍNH========================================
\color{black}

\txtbx[2]{
    {\color{vol} \uline{Tóm tắt:}} Tình yêu mãnh liệt với hóa học nhen nhóm từ thuở thơ ấu của Suzuno, từng bị chôn vùi trong những năm tháng gian khó và mất mát ở thế giới cũ, đã bùng cháy trở lại sau khi được anh trai Kuon cứu vớt. Thế nhưng, khi những phản ứng tổng hợp ngày càng phức tạp và đòi hỏi sự tinh vi hơn, góc nhỏ chật chội mà cô đang mượn ở tầng bốn bỗng trở thành một giới hạn ngột ngạt và nguy hiểm. Khao khát về một không gian nghiên cứu độc lập, an toàn cho cả gia đình đã thôi thúc cô bé đi đến một quyết định táo bạo: tìm kiếm một phòng thí nghiệm thực thụ cho riêng mình.
}

\pov{Hikawa Suzuno}

Nếu có ai đó hỏi tôi rằng, khoảnh khắc nào trong đời khiến tôi nhận ra thế giới này kỳ diệu đến nhường nào, tôi sẽ không ngần ngại kể cho họ nghe về buổi chiều mùa thu năm tôi lên bảy tuổi.

Thuở ấy, khi bố mẹ tôi vẫn còn hiện diện trên cõi đời này, tôi từng là một cô bé tò mò đến mức kỳ lạ đối với mọi hiện tượng của tự nhiên. Trong khi những đứa trẻ đồng trang lứa mải mê với búp bê vải hay đuổi bắt bướm ngoài đồng, tôi lại thích ngồi hàng giờ bên bậc cửa sau của khu nhà cũ, ngắm nhìn giọt sương đọng trên phiến lá khoai nước, hay mê mẩn quan sát cách những đốm vôi sống sủi bọt ùng ục, tỏa hơi nóng ngùn ngụt khi ông nội tôi thả chúng vào vại nước tôi vôi.

Bố tôi - vốn là một người thợ cơ khí quanh năm gắn bó với dầu mỡ và máy tiện - đã mỉm cười xoa đầu tôi và bảo: ``Đó là Hóa học đấy, Suzu à. Mọi thứ trên thế giới này, từ hòn đất, ngọn cỏ, dòng nước con uống cho đến không khí con hít thở, đều được tạo nên từ những hạt li ti không thể nhìn thấy bằng mắt thường. Và khi chúng tìm thấy nhau, chúng sẽ nhảy múa, biến đổi để sinh ra những điều kỳ diệu.''

Lời nói giản dị ấy của bố đã gieo vào tâm hồn non nớt của tôi một hạt mầm đam mê cháy bỏng. Với tôi, hóa học chưa bao giờ là những công thức khô khan in trên trang sách hay những bảng tuần hoàn phức tạp. Nó là một thứ ngôn ngữ bí mật của tự nhiên, là cây đũa phép thuật của một phù thủy có thể biến những điều tầm thường nhất thành phép màu lung linh sắc màu.

Thế nhưng, giấc mơ êm đềm ấy đã tan vỡ như bọt xà phòng chỉ vài năm sau đó.

Cơn bão nghiệt ngã của số phận ập xuống gia đình tôi ở vũ trụ cũ. Mẹ tôi qua đời trong một vụ tai nạn giao thông vào một chiều mưa trơn trượt. Ba năm sau, đến lượt bố ra đi vĩnh viễn sau sự cố máy ép kim loại trong nhà máy chế tạo. Chỗ dựa duy nhất sụp đổ, gia tộc họ hàng bên nội quay lưng vì những tranh chấp đất đai cay đắng, bỏ mặc hai chị em tôi trơ trọi giữa cuộc đời.

Năm ấy, tôi vừa tròn tám tuổi, còn chị Kaiho mới mười tám.

Để kiếm từng đồng tiền còm cõi nuôi tôi ăn học và giữ cho hai đứa không bị lùa vào trại trẻ mồ côi, chị Kaiho đã phải cắn răng nghỉ học, quần quật làm thuê từ sáng sớm đến đêm muộn. Hai chị em tôi dọn về một căn phòng trọ tồi tàn chưa đầy mười lăm mét vuông ở tầng ba của một khu chung cư mini xuống cấp. Nơi ấy, cầu thang lúc nào cũng nhớp nháp nước rác, tường mốc meo nứt nẻ, và mỗi mùa mưa bão nước lại dột lênh láng qua khe trần.

Những ngày tháng ấy, bữa cơm của hai chị em chỉ là những gói mì tôm sống bẻ đôi, hay những bát cơm nguội xin từ hàng xóm ăn cùng nước tương nhạt thếch. Học phí của tôi luôn đóng muộn nhất lớp, và thậm chí tôi đã phải gián đoạn việc học tới hai lần, khiến tuổi học của tôi bị chậm hẳn hai năm so với bạn bè cùng trang lứa.

Trong một góc sống chật vật, ngột ngạt và cùng cực đến như thế, làm sao một đứa trẻ như tôi dám mở miệng xin tiền mua một ống nghiệm thủy tinh hay vài giọt dung dịch chỉ thị màu?

Ước mơ về một chiếc áo choàng blouse trắng tinh khôi, về những bình tam giác chứa dung dịch đổi màu lấp lánh đành phải lặng lẽ khép lại. Tôi giấu chặt cuốn vở ghi chép những phản ứng hóa học tự mày mò vào đáy chiếc vali sờn rách, khóa kín niềm đam mê lại như một món đồ xa xỉ không có quyền được phép chạm tới. Tôi tự nhủ rằng, chỉ cần chị Kaiho không phải gục ngã vì kiệt sức, chỉ cần hai chị em tôi còn có nhau để nương tựa, thì việc từ bỏ một ước mơ con trẻ chẳng đáng là gì.

\ct

Cho đến một ngày nọ, số phận đã dẫn lối chúng tôi bước qua bờ vực của sự tuyệt vọng.

Sau trận hỏa hoạn kinh hoàng thiêu rụi khu chung cư mini cũ và cú trượt ngã định mệnh vào khoảng không vô tận của The Void, hai chị em tôi đã gặp anh Kuon.

Giữa khoảng không hư vô lạnh lẽo và hỗn loạn của đa vũ trụ, người thanh niên mang ánh mắt kiên định ấy đã dang rộng vòng tay che chở cho hai chị em tôi. Anh không chỉ kiên nhẫn lắng nghe những nỗi đau chôn giấu bao năm, không chỉ dùng chiếc áo khoác ấm áp xua tan đi cái lạnh của sự ghẻ lạnh xã hội, mà còn dẫn dắt chúng tôi vượt qua những vòng lặp ma quái để trở về với thế giới của anh, một thế giới thanh bình nơi đại gia đình Hikawa đang giang rộng vòng tay đón chào.

Được bước đi dưới bầu trời trong xanh của thủ đô Kawachi, được sống trong căn nhà ống ấm áp bên dải công viên ven sông, và trên hết, được sống bên cạnh một người anh cả luôn ân cần, đáng tin cậy như anh Kuon, tâm hồn đầy vết thương của tôi đã dần được chữa lành.

Và chính tại nơi này, ngọn lửa đam mê từng bị tro tàn phủ kín thuở ấu thơ bỗng cựa mình thức giấc, bùng cháy mãnh liệt hơn bao giờ hết.

Tôi nhớ mãi cái ngày đầu tiên bước chân vào ngôi trường mới ở Kawachi. Khi thầy giáo môn Hóa học làm thí nghiệm cho dung dịch phenolphthalein nhỏ vào cốc nước vôi trong, nhìn giọt dung dịch không màu vừa chạm mặt nước bỗng bùng lên sắc hồng cánh sen rực rỡ, trái tim tôi đã đập rộn ràng như muốn nhảy khỏi lồng ngực. Đôi bàn tay tôi run rẩy siết chặt lấy mép bàn gỗ. Cảm giác thân thuộc, niềm say mê kỳ diệu từ những lời dặn của bố năm nào bỗng cuồn cuộn chảy trong huyết quản.

Tôi bắt đầu dành dụm từng đồng tiền tiêu vặt mà anh Kuon và chị Kaiho cho mỗi tuần để mua sách chuyên khảo. Từ những kiến thức cơ bản trong sách giáo khoa, tôi đọc ngấu nghiến lên cấp ba, rồi tìm đọc các tài liệu hóa học hữu cơ, hóa sinh nâng cao của bậc đại học.

Anh Kuon nhận ra niềm yêu thích đặc biệt ấy của tôi. Anh không hề ngăn cấm hay gạt đi như những người lớn thực dụng khác, mà trái lại, anh còn mỉm cười xoa đầu tôi, đặt mua tặng tôi một bộ dụng cụ thí nghiệm cơ bản bằng thủy tinh borosilicate chịu nhiệt cao cấp nhân dịp sinh nhật tôi tròn mười sáu tuổi.

Khỏi phải nói tôi đã hạnh phúc đến nhường nào. Tôi nâng niu từng chiếc ống nghiệm, từng chiếc phễu chiết, kẹp gỗ như những báu vật vô giá.

Từ những thí nghiệm đơn giản như trích ly tinh dầu sả, kết tinh đường phèn hay thử nghiệm các chất chỉ thị màu tự nhiên từ bắp cải tím, tôi bắt đầu tiến xa hơn vào các phản ứng tổng hợp este tạo mùi hương, điều chế xà phòng handmade từ dầu dừa và sáp ong, cho đến những phản ứng trao đổi ion tinh vi hơn. Trong mắt tôi, mỗi giọt hóa chất rơi xuống, mỗi làn khói trắng bốc lên hay từng tinh thể lấp lánh kết tủa dưới đáy ống nghiệm đều là một vũ điệu tuyệt mỹ của tạo hóa. Các anh chị em trong nhà bắt đầu gọi tôi bằng cái biệt danh đáng yêu: ``Cô phù thủy hóa học''.

\ct

Thế nhưng, khi ngọn lửa đam mê càng rực cháy, thực tế nghiệt ngã lại một lần nữa gõ cửa.

Càng dấn sâu vào con đường nghiên cứu, các phản ứng của tôi càng trở nên phức tạp, đòi hỏi thiết bị gia nhiệt, hệ thống ngưng tụ hồi lưu và cả những chất phản ứng nhạy cảm hơn. Và lúc này, vấn đề nan giải nhất đã phát sinh: không gian.

Sau đợt gia đình tôi canh tân lại căn phòng trống ở tầng ba làm nơi sinh hoạt riêng cho phái nữ, phòng ngủ của tôi, chị Kaiho, Ryuuhou và bé Naomi đã được bố trí ổn định. Tuy nhiên, tầng ba là phòng ngủ chung, nơi bé Naomi hay chạy nhảy đùa nghịch và chị Kaiho cần sự ngăn nắp tuyệt đối; chỉ cần một giọt axit loãng hay một mùi hương hóa chất nồng nặc thoát ra thôi cũng đủ để biến nơi ấy thành thảm họa.

Chính vì thế, tôi đành phải chuyển toàn bộ ``gia tài'' thí nghiệm của mình lên một góc nhỏ ở tầng bốn.

Tầng bốn vốn được chia đôi: nửa bên trái là phòng khách của bác Hijaki, nửa bên phải là phòng ngủ yên tĩnh của bà nội đã lớn tuổi. Bác Hijaki tốt bụng đã dành cho tôi một góc hành lang nhỏ cạnh cửa sổ thông gió, nơi tôi có thể kê vừa một chiếc bàn gỗ cũ dài một mét hai để đặt giá ống nghiệm, chiếc đèn cồn nhỏ và vài chai lọ hóa chất đã dán nhãn cẩn thận.

Lúc đầu, góc nhỏ ấy là một chốn thiên đường đối với tôi. Cứ sau mỗi buổi tan học hoặc những chiều cuối tuần rảnh rỗi, tôi lại khoác lên mình chiếc áo blouse trắng, đeo kính bảo hộ và cẩn thận cân từng miligram hóa chất dưới ánh nắng hiền hòa rọi qua ô cửa kính.

Nhưng thiên đường ấy đang nhanh chóng biến thành một cơn ác mộng chật chội và đầy âu lo.

Thứ nhất là sự thiếu hụt không gian đến nghẹt thở. Một chiếc bàn một mét hai là quá nhỏ bé cho một người muốn thực hiện các phản ứng ngưng tụ hay chưng cất hồi lưu. Dàn ống sinh hàn, bình cầu đáy tròn, bếp cách thủy và giá đỡ chiếm trọn mặt bàn, khiến tôi thậm chí không còn đủ chỗ để đặt cuốn sổ tay ghi chép dữ liệu. Mỗi lần muốn lấy một lọ dung dịch, tôi phải lách người thật khéo léo, chỉ sợ sơ ý quệt tay áo làm rơi vỡ những bình thủy tinh mỏng manh. Chưa kể, hóa chất cần được phân loại bảo quản nghiêm ngặt: chất oxy hóa phải để riêng, chất dễ cháy phải để nơi râm mát, dung dịch axit và bazơ không được để chung ngăn tủ. Góc bàn nhỏ bé của tôi hoàn toàn bất lực trước những quy chuẩn an toàn cơ bản ấy.

Và thứ hai - điều khiến tôi trằn trọc mất ngủ suốt nhiều đêm qua - chính là mối nguy hiểm tiềm tàng đối với các thành viên cùng tầng.

Tầng bốn là nơi nghỉ ngơi của bà nội. Bà tuổi đã cao, sức khỏe yếu, hệ hô hấp rất nhạy cảm với những mùi lạ. Dù tôi đã cố gắng chọn những phản ứng ít sinh khí và luôn mở toang cửa sổ thông gió mỗi khi thao tác, nhưng những phản ứng nhiệt nhôm hay este hóa vẫn không tránh khỏi việc phát tán những phân tử mùi cay nồng thoang thoảng qua khe cửa.

Mới tuần trước thôi, khi tôi đang đun nóng một hỗn hợp axit axetic và rượu etylic để điều chế etyl axetat, một cơn gió ngược từ dải công viên ven sông bất ngờ thốc mạnh qua cửa sổ. Mùi giấm chua gắt và hơi cồn bốc ngược vào hành lang, khiến bà nội ở phòng đối diện ho sặc sụa suốt một hồi dài. Bác Hijaki phải vội vàng dìu bà ra ban công hít thở không khí trong lành.

Nhìn cảnh tượng ấy, nhìn gương mặt nhăn nheo vì mệt mỏi của bà nội và ánh mắt lo lắng của bác Hijaki, ngực tôi thắt lại vì ân hận và tội lỗi. Tôi vội vàng dập tắt đèn cồn, lau dọn sạch sẽ rồi cúi đầu lí nhí xin lỗi bà và bác. Dù bà nội chỉ hiền hậu xoa đầu tôi, bảo rằng ``Không sao đâu cháu gái, bà chỉ bị hắt hơi chút thôi'', nhưng cảm giác cắn rứt trong lòng tôi không thể nào vơi bớt.

Chưa hết, còn có cả bé út Naomi. Con bé rất thích tôi, cứ mỗi lần thấy tôi lên tầng bốn là lại lẽo đẽo chạy theo, kiễng chân tròn xoe mắt nhìn những chiếc bình thủy tinh đựng dung dịch màu xanh lam của đồng sunfat hay màu tím biếc của thuốc tím kali pemanganat. Trẻ con vốn hiếu động và ngây thơ; chỉ cần tôi lơ là một tích tắc, nếu con bé vô tình chạm tay vào hóa chất ăn mòn hay làm đổ giá ống nghiệm, hậu quả sẽ khủng khiếp đến mức nào?

Đó là một mối nguy hiểm hiện hữu từng ngày, từng giờ.

Tôi đứng tựa lưng vào lan can ban công tầng bốn, đưa mắt nhìn xa xăm về phía bờ sông rợp bóng cây xanh. Làn gió hạ mang theo hơi nước mát rượi mơn man trên má tôi, nhưng trong lòng tôi lại là một mớ bòng bong trăn trở.

Tôi không thể dừng niềm đam mê này lại được nữa. Hóa học đã hòa vào từng nhịp thở, từng giấc mơ của tôi, là con đường mà tôi muốn gắn bó trọn đời để sau này có thể tạo ra những dược liệu, những chế phẩm sinh học giúp ích cho đời và đền đáp công ơn nuôi dưỡng của anh Kuon, của bố mẹ và chị Kaiho.

Nhưng tôi cũng tuyệt đối không thể tiếp tục biến góc hành lang tầng bốn này thành một mối đe dọa rình rập đối với sức khỏe của bà nội và sự an toàn của cả gia đình.

Tôi cần một không gian rộng lớn hơn.

Một không gian thực thụ, biệt lập, có hệ thống thông gió đạt chuẩn, có khoảng cách ly an toàn với khu vực sinh hoạt của mọi người, nơi tôi có thể thỏa sức bày biện dàn thiết bị chưng cất, nơi tôi có thể đắm chìm vào những phản ứng thăng hoa mà không phải nơm nớp lo sợ làm kinh động đến bất kỳ ai.

Nhưng giữa lòng một khu nhà ống đô thị tấc đất tấc vàng như Kawachi, tìm đâu ra một chốn dung thân lý tưởng như thế cho một cô bé mười sáu tuổi?

Tôi cắn nhẹ môi dưới, ánh mắt dị sắc - một bên đỏ thẫm của anh Kuon, một bên xanh lam của chị Kaiho - bỗng ánh lên một tia quyết tâm kiên định.

Tôi phải nói chuyện này với anh Kuon. Chỉ có người anh cả thông tuệ và bản lĩnh ấy mới có thể giúp tôi tìm ra lối thoát cho bài toán hóc búa này.

\ct

Bữa cơm tối hôm đó ở gian bếp tầng năm diễn ra trong bầu không khí ấm cúng quen thuộc. Chờ cho đến khi mọi người đã dùng xong bữa, Naomi ngoan ngoãn mang đĩa hoa quả gọt sẵn đặt lên bàn trà, tôi mới hít một hơi thật sâu, ngập ngừng đem toàn bộ những trăn trở của mình bày tỏ trước cả nhà.

Nghe tôi giãi bày xong về việc góc bàn nhỏ ở tầng bốn không còn đủ chỗ và nguy cơ khói khí độc có thể ảnh hưởng đến sức khỏe của bà nội, cả nhà rơi vào một khoảng lặng suy ngẫm.

Anh Kuon đặt tách trà sen xuống khay gỗ, gật đầu vẻ đồng tình: ``Anh hiểu nỗi lòng của Suzuno. Niềm đam mê nghiên cứu của em là điều rất đáng quý, và cả nhà luôn ủng hộ hết mình. Nhưng đúng là việc để các hóa chất phản ứng ngay sát phòng ngủ của bà nội và không gian sinh hoạt chung là quá rủi ro. Chúng ta phải tìm cho em một không gian riêng biệt.''

``Hay là đưa lên sân thượng nhà mình đi anh?'' Ryuuhou vừa nhai miếng táo giòn rụm vừa hào hứng hiến kế đầu tiên. ``Trên đó thoáng gió lồng lộng, có làm nổ tung cái gì thì khói cũng bay thẳng lên trời, chẳng sợ ai ngửi thấy!''

Tôi vội vàng xua xua tay, khẽ lắc đầu phân tích: ``Không được đâu Ryuuhou-chan à. Trên đó có khu vườn nhỏ trồng rau của mẹ và giàn phơi quần áo của cả nhà. Khí clo, lưu huỳnh đioxit hay hơi axit bốc lên gặp sương đêm sẽ tạo thành mưa axit cục bộ, làm chết rụi hết cây cối và ăn mòn mục nát quần áo ngay. Chưa kể, nếu có sự cố tràn đổ hóa chất lỏng, nó sẽ ngấm thẳng xuống trần tầng dưới, cực kỳ nguy hiểm.''

Ryuuhou gãi đầu cười trừ, lè lưỡi: ``Nghe ghê quá\dots\ Thế ở phòng thí nghiệm của trường học thì sao chị? Trường cấp hai rộng thênh thang, thiếu gì phòng thực hành đâu?''

Nghe đến đây, tôi hơi đỏ mặt, cúi gằm xuống mân mê vạt áo, lí nhí đáp: ``Ở trường\dots\ phòng thực hành chỉ mở cửa đúng giờ tiết học thôi. Vả lại\dots\ em cũng bị các thầy cô trong tổ Tự nhiên cấm tiệt không cho bén mảng tới gần phòng kho hóa chất ngoài giờ rồi\dots''

Chị Kaiho đang ngồi tựa lưng vào mép ghế, nhướng mày tò mò: ``Sao lại bị cấm? Em ngoan ngoãn thế kia mà quậy phá gì ở trường à?''

Tôi ngượng ngùng gật đầu thú nhận: ``Tháng trước\dots\ vì muốn kiểm chứng tốc độ phản ứng oxi hóa khử của hợp chất peoxit ở nồng độ cao mà phòng thí nghiệm không có giờ thực hành tự do, em đã lén mượn chìa khóa của bác lao công vào phòng thí nghiệm lúc năm giờ chiều. Kết quả là\dots\ áp suất bình kín tăng quá nhanh, bọt khí màu tím trào ngược ra ngoài, làm cháy đen một góc rèm cửa và nhuộm hồng luôn cả bức tường thạch cao của trường\dots\ Em bị ghi sổ đầu bài và thầy hiệu phó bắt viết bản kiểm điểm sâu sắc, cấm cửa luôn từ đợt đó ạ\dots''

Cả gian phòng tầng năm chết lặng trong ba giây. Ryuuhou ngồi cạnh tôi, mắt tròn xoe, mồ hôi hột túa ra ròng ròng trên trán. Cô nhóc tomboy ngổ ngáo thường ngày vốn nổi tiếng bạt mạng mà giờ đây cũng phải nuốt nước bọt đánh *ực* một cái, thốt lên đầy kinh hãi: ``Trời đất quỷ thần ơi\dots\ Bình thường nhìn chị Suzu hiền như cục đất, ai ngờ lại là trùm khủng bố ngầm! Đam mê tới mức đột nhập phòng thí nghiệm trường làm nổ tường luôn kìa\dots''

Chị Kaiho cũng phải bật cười lắc đầu, còn anh Kuon thì xoa xoa sống mũi với vẻ mặt dở khóc dở cười. Anh đưa mắt nhìn sang chiếc đồng hồ thông minh trên cổ tay: ``Thế còn mang sang văn phòng \hll\ hoặc căn mancave của mấy anh \hls\ thì sao?''

Chị Kaiho lập tức giơ tay phản bác ngay tắp lự: ``Tuyệt đối không được! Bên hololive người ra vào nườm nượp, toàn mấy cô nàng idol ngáo ngơ hậu đậu, để hóa chất ở đấy khéo có người tưởng nhầm si-rô ngọt mà uống vào thì toi mạng. Còn bên mancave của tụi con trai holostars thì bừa bộn như bãi chiến trường, bọn nó quậy phá hò hét suốt ngày, chỉ cần đá bóng trúng giá kính của Suzu một cái là tan tành hết.''

Các phương án lần lượt bị gạt bỏ, bài toán không gian dường như rơi vào ngõ cụt.

Căn phòng lại chìm vào im lặng. Tôi khẽ thở dài, trong lòng dâng lên một nỗi hụt hẫng nhè nhẹ. Có lẽ ước muốn có một phòng thí nghiệm riêng ở độ tuổi này quả thực là quá xa xỉ và viển vông.

Nhưng đúng lúc ấy, chị Kaiho bỗng khựng lại, đôi mắt màu xanh lam sắc sảo khẽ chớp lên một tia sáng: ``Khoan đã\dots\ Anh Kuon, anh có nhớ bãi đất hoang nằm ngay sát vách xưởng cơ khí của gia đình mình ở ngoại ô không?''

Anh Kuon hơi nhíu mày, lục lại trí nhớ: ``Bãi đất ở ngoại ô à? Ý em là chỗ nhà máy cũ ấy hả?''

``Đúng rồi!'' chị vỗ đùi một cái rõ to. ``Ngay cạnh xưởng tiện của bố có một khu nhà xưởng bỏ hoang. Nguyên là xưởng sản xuất đá lát vỉa hè và gạch men thủ công, nhưng đã đóng cửa cách đây không lâu do làm ăn thua lỗ. Hồi nhà máy đó tuyên bố phá sản, bố Ryujin từng bỏ tiền mua lại toàn bộ khu đất ấy. Bố bảo định bụng sau này mở rộng quy mô xưởng cơ khí làm kho chứa phôi thép và máy móc cồng kềnh, nhưng đợt này công việc bận rộn quá nên bố vẫn để trống, chưa kịp sửa sang gì cả!''

Đôi mắt tôi bỗng sáng rực lên như bắt được vàng. Một bãi đất nằm ngay cạnh xưởng cơ khí của bố ở vùng ngoại ô! Nơi đó cách biệt hoàn toàn với khu dân cư đông đúc, không khí thoáng đãng, lại có bố và chị Kaiho thường xuyên làm việc bên cạnh để trông coi!

Anh Kuon khoanh tay gật gù, khóe môi khẽ cong lên: ``Ý kiến hay đấy. Ngày mai là thứ Bảy, anh với Kaiho cũng có lịch xuống xưởng phụ bố bảo dưỡng máy tiện kim loại. Tiện thể cả nhà mình cùng đi khảo sát thực địa luôn xem sao.''

Bé út Naomi nghe thấy thế liền vỗ tay reo mừng: ``Naomi cũng muốn đi xem xưởng đá lát với chị Suzuno ạ!''

Thế là, một tia hy vọng rực rỡ bất ngờ thắp sáng màn đêm trăn trở của tôi.

\ct

Sáng thứ Bảy hôm sau, cả gia đình chúng tôi rời trung tâm thành phố Kawachi trên chuyến xe buýt sớm hướng về vùng ngoại ô thanh bình.

Băng qua những dải đồng cỏ xanh mướt và những rặng tre rì rào trong gió sớm, chiếc xe dừng lại ở con đường dẫn vào xưởng cơ khí của bố. Bãi đất hoang của nhà máy đá lát cũ nằm lọt thỏm ngay bên cạnh, chỉ ngăn cách với xưởng của bố bằng một bờ rào dây thép gai đã hơi rỉ sét.

Đó là một khu nhà xưởng một tầng rộng chừng hơn sáu mươi mét vuông. Mặc dù đã bị bỏ hoang một thời gian, bụi bặm và lá khô phủ đầy lối đi, nhưng kết cấu tổng thể lại vững chãi đến bất ngờ. Bộ khung thép tiền chế sơn chống rỉ màu xám đậm vẫn còn nguyên vẹn, trần tôn cách nhiệt cao ráo, thoáng khí, và đặc biệt là toàn bộ mặt sàn đều được đổ bê-tông cốt thép dày dặn, phẳng lì; đó là những dấu tích của một nơi từng chịu tải trọng của hàng chục tấn máy ép thủy lực và đá tảng.

Vừa bước qua cánh cửa sắt lớn hoen rỉ, một làn gió mát từ cánh đồng thốc qua những ô cửa sổ thông gió trên cao khiến không gian bên trong vô cùng thông thoáng, không hề có cảm giác ẩm mốc, bí bách.

``Kết cấu tuyệt vời hơn anh nghĩ đấy,'' anh Kuon ngước nhìn trần xưởng, dùng mũi giày gõ gõ xuống nền bê-tông đặc quánh. ``Nền móng kiên cố thế này thì tha hồ kê bàn đá chống rung để đặt cân phân tích điện tử. Trần cao, cửa sổ đối lưu hai đầu thế này thì việc lắp đặt chụp hút khói và ống xả khí thải hóa chất sẽ cực kỳ thuận lợi.''

Chị Kaiho khoanh tay, đưa mắt đo đạc cự ly theo thói quen của một người thợ cơ khí lành nghề: ``Hệ thống đường cấp thoát nước công nghiệp từ hồi sản xuất đá mài vẫn còn đường ống chờ ở góc tường kia kìa. Chỉ cần xúc rửa lại và lắp thêm bể lắng lọc trung hòa axit là đạt chuẩn môi trường ngay. Bố mua lại chỗ này đúng là món hời thật sự.''

Trong khi anh Kuon và chị Kaiho mải mê đánh giá kỹ thuật, trái tim tôi đã đập rộn ràng như mở hội. Không gian này\dots\ diện tích này\dots\ nó vượt xa mọi giấc mơ hoang đường nhất mà tôi từng dám nghĩ tới!

Không thể kìm nén được sự hào hứng đang dâng trào tột độ, tôi vội vàng nhặt một mẩu phấn trắng rơi vãi bên chân tường gỉ sét, cúi người thoăn thoắt vẽ ngay trên nền bê-tông xám xịt.

Từng đường nét dứt khoát, rành mạch nhanh chóng hiện lên dưới bàn tay tôi. Khu vực bàn đá trung tâm đặt dàn chưng cất hồi lưu; góc phía đông có cửa sổ lớn làm tủ hút khí độc chuyên dụng; góc phía bắc xây tủ bảo quản hóa chất ba ngăn cách ly; góc phía nam đặt kệ để bình chữa cháy CO2 và vòi sen tắm khẩn cấp rửa mắt; lối đi giữa các phân khu rộng rãi đạt chuẩn an toàn phòng chống cháy nổ\dots

Tôi vừa vẽ vừa lẩm bẩm tính toán lưu lượng khí và độ dốc thoát nước, hai mắt lấp lánh như một đứa trẻ vừa bước vào xứ sở thần tiên.

Ryuuhou bước lại gần, cúi đầu nhìn những vòng tròn, mũi tên chằng chịt và những khối hộp vuông vức tôi vừa phác thảo trên mặt sàn, hai mắt con bé hoa lên như nhìn vào ma trận bát quái: ``Khiếp thế! Chị Suzu ơi, chị đang vẽ cái gì trông đồ sộ như boong-ke quân sự với pháo đài phòng thủ thời trung cổ thế này?!''

Tôi ngẩng lên, quệt vệt bụi phấn dính trên chóp mũi, nở một nụ cười rạng rỡ và đầy tự hào: ``Pháo đài gì chứ! Đây là sơ đồ bố trí một phòng thí nghiệm hóa học tiêu chuẩn cấp quốc tế đấy!''

Ở góc xưởng đối diện, bé út Naomi đang ngồi xổm bên một chiếc thùng gỗ mục nát cũ kỹ. Đôi mắt màu xanh lục bảo của con bé mở to, ánh lên sự mê mẩn thuần khiết khi nhìn vào những viên đá lát mẫu bị bỏ quên dưới đáy thùng.

Đó là những phiến đá hoa cương và cẩm thạch nhân tạo đủ sắc màu: có viên màu xanh ngọc bích lấp lánh vân mây trắng, có viên lại ánh lên sắc hồng ruby hay lấm tấm những hạt kim tuyến thủy tinh khúc xạ ánh mặt trời rực rỡ. Đôi bàn tay nhỏ nhắn của Naomi cẩn thận nhặt từng viên đá bóng loáng lên lau chùi, nâng niu như những viên ngọc quý giữa đống hoang tàn: ``Đẹp quá\dots\ Chị Suzu ơi nhìn này, những hòn đá ở đây phát sáng đẹp như những ngôi sao nhỏ vậy! Phòng thí nghiệm mới của chị Suzu sau này nhất định sẽ lấp lánh như một cung điện đá quý cho mà xem!''

Tiếng nói cười trong trẻo, hồn nhiên của Naomi vang vọng khắp khu nhà xưởng hoang vu, xua tan đi sự lạnh lẽo của nhiều năm tháng bị lãng quên.

Tôi đứng giữa bãi đất hoang ngập tràn ánh nắng ban mai, hít một hơi thật sâu lồng ngực mùi gió đồng nội ngai ngái hòa lẫn hương gỉ sắt cũ kỹ. Một cảm giác xúc động nghẹn ngào khó tả dâng lên nơi khóe mắt. Nơi đây, chính trên mảnh đất hoang sơ này, ước mơ cháy bỏng suốt bao năm tháng thăng trầm của tôi sắp sửa được đâm chồi nảy lộc.

\ct

Ngay chiều hôm ấy, sau khi trở về từ xưởng, anh Kuon đã trực tiếp gọi điện cho bố Ryujin để xin phép sử dụng khu nhà xưởng bỏ hoang.

Đầu dây bên kia, bố vừa nghe anh Kuon trình bày nguyện vọng của tôi thì bật cười hào sảng: ``Tưởng chuyện gì to tát, chứ Suzu muốn dùng chỗ đó để học hành nghiên cứu thì bố đồng ý cả hai tay! Nhưng giao kèo trước nhé: Bố là thợ cơ khí cả đời, cái gì dính đến sắt thép kết cấu là phải chắc như đinh đóng cột. Để đầu tuần bố cùng mấy chú thợ già qua thẩm định lại toàn bộ móng cọc và xà gồ thép chịu lực, bố duyệt xong mới được bắt tay vào làm đấy!''

Và không chỉ dừng lại ở sự hỗ trợ của bố, anh Kuon còn kích hoạt một ``lực lượng viện binh'' vô cùng đặc biệt.

Sáng thứ Ba tuần kế tiếp, một chiếc xe bán tải chở đầy thiết bị đo trắc địa laser, máy siêu âm cốt thép và máy toàn đạc điện tử bất ngờ đỗ xịch trước cổng bãi đất hoang. Bước xuống xe là năm anh sinh viên quen thuộc đến từ Khoa Thiết kế Công trình - những người bạn học xuất sắc từng giúp anh Kuon và bố khảo sát, chỉnh sửa và mở rộng căn phòng gác mái ở ký túc xá hồi trước\footnote{{\flaregothic HISTOIRE 3}, quyển Ký túc xá.}.

``Chào nhóc phù thủy nhỏ!'' anh chàng đội trưởng nhóm sinh viên bước tới, giơ tay vẫy chào tôi đầy niềm nở. ``Nghe Kuon bảo em gái cậu ấy chuẩn bị xây một trung tâm nghiên cứu ma thuật bí mật nên bọn anh phải tức tốc phóng xe xuống đây ngay. Yên tâm, với tay nghề của bọn anh thì dù em có muốn biến cái xưởng gạch cũ này thành viện hàn lâm khoa học cũng xong tuốt!''

Cả khu xưởng hoang phế bỗng chốc biến thành một công trường sôi động hơn bao giờ hết.

Dưới sự chỉ đạo của bố Ryujin và sự tính toán chuẩn xác từ nhóm sinh viên kiến trúc, từng mảng việc được phân chia nhịp nhàng như một cỗ máy vận hành trơn tru.

Phần việc nặng nhọc và đòi hỏi sức lực cơ bắp nhất đương nhiên được giao cho chị Kaiho. Gỡ bỏ những khối bệ máy ép thủy lực rỉ sét nặng hàng tấn hay di dời những tảng đá hoa cương phế liệu nằm chắn lối đi, thì với người bình thường thì phải dùng đến xe cẩu chuyên dụng, nhưng đối với chị ấy, đó lại là cơ hội để cô em gái song sinh của anh Kuon thể hiện sức mạnh thể chất phi thường. Chị tháo chiếc nơ kẹp tóc hãm lực quen thuộc, xắn tay áo bảo hộ rồi nhẹ nhàng nhấc bổng từng khối kim loại nặng cả ngàn cân ném gọn vào thùng xe phế liệu như thể đang dọn dẹp mấy hộp các-tông rỗng, khiến mấy anh sinh viên kiến trúc đứng nhìn trân mà toát cả mồ hôi hột.

Về phần mình, tôi phụ trách trái tim vận hành của phòng thí nghiệm: thiết kế hệ thống thông gió và xử lý nước thải hóa chất.

Tôi không muốn bất kỳ giọt hóa chất độc hại nào thấm vào lòng đất hay mạch nước ngầm của vùng ngoại ô này. Tôi ngồi bệt xuống một góc sàn sạch sẽ, miệt mài vẽ sơ đồ đường ống dẫn khí thải từ chụp hút trung tâm qua tháp hấp thụ dung dịch kiềm để trung hòa hơi axit trước khi xả ra ngoài không khí. Bên cạnh là hệ thống bể lắng ba ngăn lọc than hoạt tính và sỏi thạch anh để lọc sạch nước rửa dụng cụ.

Mỗi khi tôi vướng mắc về áp suất luồng khí hay độ chịu tải của hệ thống quạt hút, từ chiếc đồng hồ thông minh trên cổ tay anh Kuon, giọng nói điềm tĩnh của Game lại cất lên: ``Thể tích xưởng tương đương một trăm tám mươi mét khối. Với tần suất trao đổi khí tiêu chuẩn cho phòng thí nghiệm cấp độ hai, quạt hút hướng trục công suất nửa kilowatt là chưa đủ áp suất bù trừ tĩnh. Tôi nghĩ cô nên nâng cấp lên quạt ly tâm công suất một phẩy hai kilowatt, bố trí van một chiều chống chảy ngược để ngăn hiện tượng đoản mạch dòng khí.''

``Ông ấy tính toán chuẩn xác đấy,'' anh chàng sinh viên kiến trúc gật gù thán phục khi nhìn vào biểu đồ mô phỏng 3D do Game phóng chiếu. ``Để bọn anh tính lại độ võng của dầm thép đỡ quạt trên mái tôn. Đảm bảo triệt tiêu toàn bộ rung chấn khi động cơ chạy hết công suất!''

Thậm chí, ngay cả Ryuuhou và bé Naomi cũng không hề chịu đứng ngoài cuộc. Ryuuhou phụ trách việc cầm chổi đót và vòi xịt nước áp lực cao quét sạch bụi bặm nhiều năm trên từng vách tường, thi thoảng lại cất giọng hát vài câu nhạc rock khàn khàn cho cả công trường thêm phần náo nhiệt. Còn Naomi thì tỉ mỉ đem những viên đá lát cẩm thạch đủ màu nhặt được từ hôm trước, dùng vữa xi măng gắn quanh bồn cây nhỏ trước cửa xưởng, tạo nên một lối vào rực rỡ sắc màu đúng như lời hứa của con bé.

Suốt gần ba tuần ròng rã, cứ sau giờ làm việc và học tập, cả nhà lại cùng nhau tề tựu dưới mái xưởng ngoại ô. Tiếng búa gõ lách cách, tiếng máy hàn xì xẹt xẹt của bố Ryujin, tiếng bàn bạc chuyên môn sôi nổi của anh Kuon và nhóm bạn học, hòa cùng mùi sơn mới thơm nồng trong gió chiều. Một tình thân gắn kết vô hình cứ thế lớn dần lên, bao bọc lấy mảnh đất này bằng tất cả tình yêu thương và sự chở che.

\ct

Một buổi chiều cuối tháng Bảy, nắng vàng như rót mật trên những cánh đồng ngoại ô Kawachi.

Tôi đứng trước cửa chính của khu nhà xưởng, hai bàn tay siết chặt lấy nhau trước ngực, lồng ngực đập rộn ràng đến mức nghẹt thở.

Trước mắt tôi, bãi đất hoang tàn u ám của một tháng trước đã biến mất hoàn toàn. Thay vào đó là một phòng thí nghiệm hóa học khang trang, hiện đại và sạch bong không một hạt bụi.

Dưới ánh đèn huỳnh quang chống nổ sáng rực, dãy bàn thí nghiệm bằng đá granite đen tuyền chống ăn mòn hóa chất chạy dài tít tắp. Trên mặt bàn, những giá đỡ bằng thép không gỉ sáng bóng nâng niu dàn ống nghiệm, bình cầu đáy tròn, ống sinh hàn xoắn ốc bằng thủy tinh trong suốt lấp lánh như pha lê. Tủ hút khí độc chuyên dụng bằng kính cường lực sừng sững ở góc tường, hệ thống vòi sen tắm khẩn cấp và bồn rửa mắt màu vàng tươi nổi bật bên cạnh chiếc bình chữa cháy bọt tuyết đạt chuẩn. Ngoài hiên cửa, những viên đá hoa cương rực rỡ do chính tay bé Naomi gắn vẫn đang phản chiếu ánh hoàng hôn lấp lánh như những ngôi sao nhỏ.

Tôi khẽ bước vào trong, ngón tay run run chạm nhẹ lên mặt bàn đá mát lạnh. Mùi cồn sát khuẩn tinh khiết, mùi kính mới và hương thơm thoang thoảng của gió đồng lùa qua khe cửa sổ thông gió khiến sống mũi tôi cay xè. Một giọt nước mắt hạnh phúc bất giác lăn dài trên má, rơi xuống mặt bàn đá vỡ tan.

Đó không phải là nước mắt của sự tủi hờn hay đau thương. Đó là giọt nước mắt của một giấc mơ ngỡ đã chôn vùi dưới đáy tro tàn, nay đã thực sự hồi sinh rạng rỡ giữa thế gian.

Tiếng bước chân quen thuộc vang lên sau lưng. Anh Kuon bước tới bên cạnh tôi, khẽ vỗ nhẹ lên vai tôi một cái đầy ấm áp: ``Từ nay về sau, đây là vương quốc của em rồi đấy, Suzuno. Thỏa sức vẫy vùng nhé. Nhưng mà\dots\ nhớ hứa với anh một điều quan trọng nhất: tuyệt đối đừng cho nổ tung căn xưởng này đấy.''

Tôi quay sang nhìn anh trai, bắt gặp ánh mắt nửa trêu chọc nửa cưng chiều của anh, bất giác đưa tay gãi đầu, bật cười ngượng ngùng: ``Dạ\dots\ em chỉ hứa là sẽ tuân thủ quy trình an toàn thôi chứ\dots\ còn nổ hay không thì\dots\ em không dám hứa chắc đâu ạ, Onii-sama\dots''

Anh Kuon giật mình thót lại, ôm trán thở dài bất lực, miệng lầm bầm: ``Biết ngay là em sẽ nói vậy mà\dots'' trong khi chị Kaiho và bố Ryujin đứng cạnh chỉ biết cười vang cả một góc trời.

\omk

Ngày hôm sau, buổi vận hành thử nghiệm đầu tiên của phòng thí nghiệm mới.

Trong lúc tôi đang cẩn thận chuẩn bị cân hóa chất cho phản ứng kết tinh muối kép, Ryuuhou vì quá tò mò trước một dàn chai lọ sặc sỡ trên kệ nên đã lén lút lẻn vào lúc tôi quay lưng đi lấy nước cất.

Nhìn thấy hai chiếc cốc thủy tinh đựng dung dịch không màu đặt cạnh nhau - một bên là dung dịch axit loãng tôi vừa pha, một bên là xà phòng lỏng trộn bột nở baking soda - cô em gái tomboy ngổ ngáo bỗng nảy ra ý định ``trổ tài phù thủy''.

``Hóa học thì có cái gì khó chứ, cứ đổ lẫn vào nhau là nó tự phản ứng thôi mà!''

Nghĩ là làm, em ấy tiện tay cầm chiếc cốc đựng xà phòng đổ ụp vào cốc dung dịch axit loãng, tiện tay bốc thêm một nắm bột phẩm màu đỏ quẳng thẳng vào bên trong rồi lấy que thủy tinh khuấy lấy khuấy để.

``N-Này, Ryuuhou-chan! Đừng có nghịch hóa chất của chị như thế--!!'' tôi chỉ kịp thốt lên can ngăn. Nhưng mọi thứ đã quá muộn.

Và kết quả đến ngay tức khắc.

*XÈÈÈO\dots\ ÙNG ỤC!!*

Khí cacbonic sinh ra từ phản ứng trung hòa kết hợp với chất tạo bọt bùng lên dữ dội như một ngọn núi lửa mini phun trào. Một lớp bọt xà phòng đặc quánh màu đỏ tươi như máu ùng ục tràn ra khỏi miệng cốc, cuồn cuộn trào qua mép bàn rồi chảy lênh láng khắp mặt sàn bê-tông bóng loáng.

Đúng lúc đó, chị Kaiho từ ngoài xưởng cơ khí bước vào định gọi cả nhà uống nước. Vừa đặt chân qua ngưỡng cửa, đế giày thể thao của chị giẫm trúng bãi bọt xà phòng trơn tuột như mỡ.

*VÈO\dots\ HUỴCH!!*

Một tiếng va chạm long trời lở đất vang lên. Người con gái sở hữu sức mạnh cơ bắp phi thường có thể nâng bổng cả tấn thép, giờ đây lại bị một bãi bọt xà phòng quật ngã ngửa ra sàn, hai chân chổng thẳng lên trời trong tư thế không thể nào thảm hại hơn.

``\textbf{Hikawa Ryuuhou!!}'' tiếng gầm thét long trời lở đất của chị tôi làm rung chuyển cả mái tôn xưởng.

Ryuuhou mặt cắt không còn một giọt máu, run như cầy sấy lùi dần vào góc tường, chiếc que thủy tinh trên tay rơi đánh *keng* xuống sàn.

Anh Kuon từ ngoài cửa lững thững bước vào, nhìn bãi chiến trường đỏ loét lênh láng dưới sàn, rồi nhìn sang cô em gái thứ tư đang co rúm người lại. Anh lắc đầu, buông một tiếng thở dài thườn thượt, cất giọng châm biếm bằng vẻ mặt bình thản đến phũ phàng: ``Anh đã bảo rồi mà\dots\ Trông mong gì từ một đứa dốt đặc cán mai các môn tự nhiên như Ryuuhou chứ. Đụng vào cái gì là phá hoại cái đó.''

Ryuuhou nhăn tít mặt mũi lại, hai má phồng to như cá nóc, vừa ấm ức vừa mếu máo cãi lại: ``Onii!! Em chỉ muốn làm một cốc sinh tố dâu sủi bọt tặng chị Suzu thôi mà! Ai ngờ bọt nó lại phun ra như vòi cứu hỏa thế kia chứ!!''

Game từ trên đồng hồ khẩy một tiếng mỉa mai khô khốc kết lại buổi thử nghiệm đầy giông bão: ``Ryuuhou-user đã kích hoạt phản ứng tạo bọt khí không kiểm soát và hạ đoán knock-out một thực thể mang sức mạnh một tấn bằng một bãi xà phòng. Khả năng cao là cô đã chính thức bị tước quyền tiếp cận phòng thí nghiệm vô thời hạn nếu không có người giám hộ.''

Giữa bãi bọt xà phòng đỏ rực và tiếng cãi cọ chành chọe inh ỏi của hai chị em, tôi đứng tựa lưng vào mép bàn đá, ôm bụng bật cười nghiêng ngả. Phòng thí nghiệm mới của tôi, xem ra từ nay về sau sẽ chẳng bao giờ thiếu đi những tiếng cười rộn rã của một gia đình.

\end{document}